\section{Explanatory conditions for finite-action decisions}
\label{sec:explanatory-propositions}

Four elementary sufficient conditions distinguish changes in prediction
error from changes in a decision. They explain possible behavior of a
finite-action predictor without establishing that the fitted models meet
these conditions. Their role is interpretive rather than a new theory
contribution.

Fix a state $h$ and multiplier $\lambda$. Both the predictor and oracle use
the same nonempty finite feasible set $\mathcal A(h)$, which includes STOP,
and the same fixed deterministic order to resolve ties. For each feasible
action, write the finite real values
\[
 v_a=R(h,a)+\lambda c_{\rm pre}(a),\qquad
 \widehat v_a=\widehat R(h,a)+\lambda c_{\rm pre}(a).
\]
The cost terms are identical and fixed in the two values. Let $a^*$ and
$\widehat a$ be their respective minimizing actions, and let
$E=\max_{a\in\mathcal A(h)}|\widehat v_a-v_a|$. Thus $E$ is also the
maximum residual-prediction error on this set. All statements below are
conditional on this shared feasible set, not assertions about how an
unknown action's availability is learned.

\begin{proposition}[Common-mode invariance]
\label{prop:common-mode-invariance}
For any finite state-dependent scalar $b(h)$, replacing every predicted
value $\widehat v_a$ by $\widehat v_a+b(h)$ leaves the selected action and
every predicted STOP-relative difference unchanged. Absolute prediction
error need not remain unchanged.
\end{proposition}
\begin{proof}
For all feasible $a,a'$, the shifted difference is
$(\widehat v_a+b(h))-(\widehat v_{a'}+b(h))
=\widehat v_a-\widehat v_{a'}$. The minimizing set and its deterministic
tie resolution therefore remain the same. Taking $a'=\operatorname{STOP}$
gives the relative-difference claim. In contrast, the absolute error becomes
$|\widehat v_a+b(h)-v_a|$, which is not generally invariant.
\end{proof}

\begin{proposition}[Strict STOP/non-STOP sign stability]
\label{prop:stop-sign-stability}
For a feasible non-STOP action $a$, define
$\Delta_{\lambda,a}=v_a-v_{\operatorname{STOP}}$ and
$\widehat\Delta_{\lambda,a}=\widehat v_a-\widehat v_{\operatorname{STOP}}$.
If
\[
 |\widehat\Delta_{\lambda,a}-\Delta_{\lambda,a}|
 <|\Delta_{\lambda,a}|,
\]
then both differences have the same strict sign. In particular, their
judgments of whether $a$ strictly improves on STOP agree.
\end{proposition}
\begin{proof}
If $\Delta_{\lambda,a}>0$, the strict bound gives
$\widehat\Delta_{\lambda,a}>
\Delta_{\lambda,a}-|\Delta_{\lambda,a}|=0$.
If $\Delta_{\lambda,a}<0$, it gives
$\widehat\Delta_{\lambda,a}<
\Delta_{\lambda,a}+|\Delta_{\lambda,a}|=0$.
The hypothesis cannot hold when $\Delta_{\lambda,a}=0$, so it makes no
claim about a true STOP tie. A negative difference denotes improvement.
\end{proof}
At $\lambda=0$ these are the residual differences used in DRRE. Correctly
classifying one action relative to STOP does not by itself identify the
best repair among all actions.

\begin{proposition}[Action identity under a positive margin]
\label{prop:action-identity}
Suppose the feasible set has at least two actions and the true minimizer
$a^*$ is unique. Define its best--second-best margin by
\[
 m=\min_{a\ne a^*}(v_a-v_{a^*})>0.
\]
If $E\leq\epsilon$ for some finite $\epsilon\geq0$ and $m>2\epsilon$,
then the predicted minimizer is uniquely $a^*$.
\end{proposition}
\begin{proof}
For every $a\ne a^*$,
\[
 \widehat v_a-\widehat v_{a^*}
 \geq (v_a-\epsilon)-(v_{a^*}+\epsilon)
 \geq m-2\epsilon>0.
\]
Thus every other predicted value is strictly larger. The non-strict
condition $m\geq2\epsilon$ alone would allow a predicted tie and would not
ensure identity under an arbitrary fixed tie order.
\end{proof}

\begin{proposition}[Margin-dependent regret]
\label{prop:margin-regret}
On a feasible set with at least two actions, use the preceding margin when
the true minimizer is unique, and set $m=0$ if it is tied. Then
\begin{equation}
 0\leq v_{\widehat a}-v_{a^*}
 \leq 2E\,\mathbf 1\{E\geq m/2\}.
 \label{eq:explanatory-margin-regret}
\end{equation}
Across a population of such decisions, if $E\leq\epsilon$ almost surely
for a fixed finite $\epsilon\geq0$, then
\begin{equation}
 \mathbb E[v_{\widehat a}-v_{a^*}]
 \leq 2\epsilon\,\Pr(m\leq2\epsilon).
 \label{eq:explanatory-expected-regret}
\end{equation}
\end{proposition}
\begin{proof}
True optimality gives nonnegative regret. Predicted optimality gives
$\widehat v_{\widehat a}\leq\widehat v_{a^*}$, hence
\[
 v_{\widehat a}-v_{a^*}
 =(v_{\widehat a}-\widehat v_{\widehat a})
  +(\widehat v_{\widehat a}-\widehat v_{a^*})
  +(\widehat v_{a^*}-v_{a^*})\leq2E.
\]
If $E<m/2$, the positive-margin proposition, applied with $\epsilon=E$,
gives $\widehat a=a^*$ and zero regret. Otherwise the $2E$ bound applies;
for $m=0$ it applies without an action-identity guarantee. This proves
Equation~\ref{eq:explanatory-margin-regret}. Under the almost-sure bound,
$\{E\geq m/2\}\subseteq\{m\leq2\epsilon\}$, so the right-hand side is
at most $2\epsilon\mathbf 1\{m\leq2\epsilon\}$. Taking expectations
proves Equation~\ref{eq:explanatory-expected-regret}; no independence of
$E$ and $m$ is required.
\end{proof}

For a single feasible action, both rules select STOP and regret is zero.
There is no non-STOP comparison or second-best action; the sign statement
is vacuous and the margin argument is unnecessary. Equivalently, setting
$m=+\infty$ on these states lets both regret bounds include them with zero
contribution. The singleton and tied-optimum cases therefore require no
assumption of a positive finite margin.

These bounds concern a uniform error across the feasible actions of an
individual decision. Small average MAE or DRRE does not establish such an
error bound, nor do these within-state results guarantee the ordering of
states for sparse dispatch. The conditions themselves provide neither a
model-specific uniform-error bound nor a population margin distribution.
The historical VH3 margin trend failed to replicate
(Supplement~\ref{sec:mechanisms-history}), so these conditions do not
establish that margin explains the observed prediction--decision gap.
